%
% TODO figure out formatting of resume, it sucks rn
%

\documentclass[9pt]{article}

\usepackage{enumitem}
\usepackage{titlesec}
\usepackage{titling}
\usepackage{mathptmx}
\usepackage{hyperref}
\usepackage{xcolor}
\usepackage[margin=0.5in]{geometry}
% \usepackage{fullpage}

% I don't want indents
\setlength{\parindent}{0in}

\pagenumbering{gobble}

% Cool link colors
\definecolor{myblue}{rgb}{0.4, 0.4, 0.8} % Define a custom blue color

\hypersetup{
    colorlinks=true,
    linkcolor=myblue,   % Internal links (sections, pages, etc.)
    citecolor=myblue,   % Citations
    urlcolor=myblue,    % URLs
}


\titleformat{\section}
{\large\bfseries}
{}
{0em}
{}[\titlerule\vspace{-0.3em}]

\titleformat{\subsection}
{\bfseries}
{}
{0em}
{\vspace{0.1em}}

\titleformat{\subsubsection}[runin]
{\bfseries}
{}
{0em}
{}[  --]

% \renewcommand{\familydefault}{\sfdefault}

\renewcommand{\maketitle}{
		\vspace{-8.5em}
    \begin{center}
		{\huge\bfseries
		\theauthor}
    \\
	    %Canadian citizen -- TN, J-1, H1-B Eligible\\
      \vspace{0.25em}
\begin{tabular}{ c | c | c }
  \href{mailto:ianayl.work@gmail.com}{\textit{ianayl.work@gmail.com}} & \href{https://www.linkedin.com/in/ianayl}{\textit{linkedin.com/in/ianayl}} & \href{https://www.github.com/ianayl}{\textit{github.com/ianayl}}
\end{tabular}
    \end{center}
    \vspace{-2.5em}
}

\titlespacing{\subsubsection}
{0em}{0.25em}{0.5em}

\titlespacing{\subsection}
{0em}{1.5em}{0em}

\setitemize{noitemsep}

\author{Ian Li}

\begin{document}

\maketitle

\section{Education}
\subsection{\textnormal{University of Western Ontario - \normalsize Honours Computer Science -- \textit{3.9/4.0 GPA}} \hfill \normalsize\textnormal{2021 - \textit{Apr. 2026}}}
\vspace{0.5em}

\section{Experience}

\subsection{Compiler Engineer Intern - \textit{Intel} \hfill \normalsize\textnormal{ May 2024 - Oct. 2025}} % \textit{18 months} --
Contributing to \textbf{MLIR} (\href{https://github.com/llvm/llvm-project}{\textit{llvm/llvm-project}}) and Intel MLIR Extensions (IMEX) for \textbf{SPIR-V} \textbf{code generation}:
\vspace{-0.5em}
\begin{itemize}
  \item Wrote MLIR's \texttt{MathToXeVM} pass, converting \textit{fastmath} ops to equivalent native GPU instructions (C++, LLVM, MLIR, SPIRV), \textbf{accelerating AI/ML, vector ops} on \textbf{Intel GPUs} 
  \item Added lowerings for MLIR's \texttt{ArithToSPIRV} pass, fixing crashes when converting \texttt{arith.index\_cast} ops to SPIRV ops 
\end{itemize}
\vspace{-0.25em}
Contributed 100+ PRs to DPC++ SYCL compiler (\href{https://github.com/intel/llvm}{\textit{intel/llvm}}), enabling GPU programming on \textbf{Intel GPUs} (Data Center GPUs (\textit{PVC, Gaudi}), Discrete GPUs (\textit{Arc, Battlemage}), \textit{Aurora Supercomputer}, Iris Xe Integrated GPUs):
\vspace{-0.5em}
\begin{itemize}
  \item Developed \textbf{performance benchmarking} CI (SYCL, Python, Shell, C++), aiding with \textbf{performance refactoring} and automatically \textbf{detecting regressions} pushed to \href{https://github.com/intel/llvm}{\textit{intel/llvm}} within 24 hours
  \item Implemented \textbf{GPU-host shared memory prefetch} for DPC++, improving performance by \textbf{eliminating latency} from \textit{page faults}
  \item Debugged critical issues in DPC++ compiler across diverse environments, ranging from deadlocks in distributed PyTorch models (\textit{Intel PyTorch Extension}) to debugging faulty GPU IR behavior (\textit{GenISA, SPIRV, IGC vISA})
  %\item Investigated instrumenting DPC++ programs with LLVM xray for cycle-level accounting of runtime performance (Clang driver)
\end{itemize}
\vspace{-1em}

%\subsection{Improving KLARAPTOR: Automatic Finding of Optimal CUDA Kernel Launch Parameters\hfill \normalsize\textnormal{Current}}
\subsection{Research Assistant - \textit{Ontario Research Centre for Computer Algebra}\hfill \normalsize\textnormal{Sep. 2023 - Apr. 2024}} % \textit{8 months} -- 
Researched methods for finding range of \textbf{CUDA} kernel launch configurations, enabling fast \textbf{autotuning} for \textbf{GPU compilers} compared to traditional autotuners while also \textit{requiring no manual user input}
\vspace{-0.5em}
\begin{itemize}
  \item Implemented \textbf{LLVM} pass for extracting dimensionality of GPU kernel from \textbf{LLVM IR} (via clang), helping to determine problem space of kernel launch parameters for GPU kernel autotuning
  \item Researched automatic deduction of \textit{valid} kernel configurations (\textit{constraints}) from GPU kernel source code, helping to enable automatic \textbf{GPU kernel autotuning} and avoid erroneous kernel configurations when autotuning
  \item Created static code analysis pass (\textbf{Clang Frontend}, Quantifier Elimination via \href{https://maplesoft.com}{\textit{Maple}}) to automatically derive GPU kernel configuration \textit{constraints} from \textbf{CUDA} source code
  % \item Working on rewrite using a \textbf{LR} parser with integrated lexer, a compiler to \textbf{LLVM} IR, and a strong typing system
        % TODO phrasing here sucks redo it
\end{itemize}
%\vspace{-0.5em}
%Goal is a novel GPU autotuning technique requiring \textit{no user intervention} while still being \textit{faster} than traditional exhaustive search
\vspace{-1em}

%\section*{Work Experience}

% \subsection{Software Developer, \textit{FIRST Robotics Team 6866} - \normalsize\textit{Markham, Canada}\hfill \normalsize\textnormal{2020}}
% \begin{itemize}
%   \item Developed a control system and API for a conveyor magazine system to shoot at targets in \textbf{Java} using WPILib
%   \item Contributed to custom PID controller class for precise control of motors
% \end{itemize}

\subsection{Lead Software Developer - \textit{Rentura (Startup)} \hfill \normalsize\textnormal{May 2021 - Feb. 2022}} % \textit{10 months} -- 
Lead the development and design of a \textit{minimum viable product} (MVP) for a B2C furniture rental platform
\vspace{-0.5em}
\begin{itemize}
	\item Designed and developed server backend (\textit{Javascript}), with a RESTful API using CRUD endpoints secured with tokens
  \item Maintained \textit{MongoDB} instances on a \textbf{Linux} server, and designed database schemas for an \textit{order management system}
  \item Developed responsive, mobile-first e-commerce frontend in \textit{React.js} using \textit{Sass} and \textit{Axios} (\textit{MERN} stack)
    % \item TODO add business stuff later
    % \item 2021\--Current
\end{itemize}
\vspace{-0.5em}

%     \item Created custom game mods and 3D assets for game servers using Unity and Blender
%     \item Maintained servers through Telnet, maintained internal management dashboards and image assets
% \end{itemize}

% \subsection{Computer Repairs} 
% \subsubsection{Home} Repaired laptop components, reflashed BIOS, fixed and reinstalled Operating Systems for clients privately 

% \subsection{Teaching Assistant}
% Ran Google CS/STEM workshops (hour of code, scratch, basic 3D modelling, etc) --- 2014\--2018  

\section{Projects}

\subsection{Experimental ML Compiler \normalsize\textnormal{- \href{https://github.com/ianayl/elohssa}{\textit{ianayl/elohssa}}}\hfill \normalsize\textnormal{May 2024 - \textit{Now}}}
%Shells are \textit{outdated}: Creating a language in \textbf{C++} with goals to be a good general-purpose language \textit{and} a fast replacement to shells
%\vspace{-0.5em}
\begin{itemize}
  \item Writing Python bindings (\textbf{pybind11}) for ML/tensor operations that can \textit{JIT}-compile into \textit{optimized GPU kernels}, akin to \textit{PyTorch}
  \item Devising \textbf{MLIR} passes for progressively lowering high-level IR to MLIR \texttt{GPU} dialect to generate/serialize GPU binaries/IR
  % \item Working on rewrite using a \textbf{LR} parser with integrated lexer, a compiler to \textbf{LLVM} IR, and a strong typing system
        % TODO phrasing here sucks balls redo it
  %\item Working on a static type-checker over the AST, and enabling JIT compilation using LLVM
\end{itemize}
\vspace{-1em}

\subsection{Experimental Interpreted Language \normalsize\textnormal{- \href{https://github.com/ianayl/ex-interpreter}{\textit{ianayl/ex-interpreter}}}\hfill \normalsize\textnormal{Jun. 2020 - Sep. 2021}}
%Shells are \textit{outdated}: Creating a language in \textbf{C++} with goals to be a good general-purpose language \textit{and} a fast replacement to shells
%\vspace{-0.5em}
\begin{itemize}
  \item Devised an automatic table-based \textbf{lexer generator}, and wrote a \textbf{recursive descent} parser following LL grammar
  \item Wrote an interpreter using a \textbf{tree-walk evaluator}, and AST passes to generate CFGs/DAG for IR generation
  % \item Working on rewrite using a \textbf{LR} parser with integrated lexer, a compiler to \textbf{LLVM} IR, and a strong typing system
        % TODO phrasing here sucks balls redo it
  %\item Working on a static type-checker over the AST, and enabling JIT compilation using LLVM
\end{itemize}
\vspace{-1em}

% \subsection{Rentura \textnormal{(Startup)} \hfill \normalsize\textnormal{2021 -- 2022}}
% Lead the development and design of a \textit{minimum viable product} for a B2C furniture rental platform
% \vspace{-0.5em}
% \begin{itemize}
%   \item Independently developed server backend in \textbf{Javascript}, with a \textbf{RESTful API} using \textbf{CRUD} endpoints secured with \textit{tokens}
%   \item Maintained \textbf{MongoDB} instances on a \textbf{Linux} server, and designed database schemas for an \textit{order management system}
%   \item Collaboratively developed responsive, mobile-first \textit{e-commerce} frontends in \textbf{React.js} using \textbf{Sass} and \textbf{Axios} (\textbf{MERN} stack)
%     % \item TODO add business stuff later
%     % \item 2021\--Current
% \end{itemize}

\subsection{Mindless \normalsize\textnormal{- \href{https://github.com/ianayl/mindless}{\textit{ianayl/mindless}}}\hfill \normalsize\textnormal{\textit{Hackwestern 2021}}}
    Proof-of-concept \textit{unconditionally secure} biometrics password manager that \textit{generates passwords} on-the-fly instead of storing them
    \vspace{-0.5em}
\begin{itemize}
	\item Devised a method to generate consistent passwords using hashing (SHA256) from \textit{facial landmark data} (using \textbf{OpenCV})
  \item Developed a password manager (\textbf{Python}, \textbf{Flask}) that \textit{generates} retrievable passwords using a \textit{face scan}
  %\item Became a finalist project in a 36-hour hackathon (Hackwestern 8) as a one-man team out of 346 teams
\end{itemize}
\vspace{-1em}

\subsection{Shell Site Generator \normalsize\textnormal{- \href{https://github.com/ianayl/shsg}{\textit{ianayl/shsg}}}\hfill \normalsize\textnormal{May 2020 - Oct. 2020}}
  Ultra lightweight and portable static site generator that will run on \textit{anything} with a minimal UNIX-like environment
    % Easy to use static site generator that can be deployed on any minimal UNIX-like environment
\vspace{-0.5em}
\begin{itemize}
  \item Developing a static site generator using POSIX shell and PCRE Regex: \textit{No dependencies other than a UNIX shell \& coreutils}
  \item Created a transpiler (Shell) for transpiling markdown to themed webpages with no alternative runtimes required
    %\item Working on removing dependency on PCRE extensions (support for Busybox), and parsing tables and lists in pure regex
\end{itemize}
\vspace{-1em}

% \subsubsection{Awards} Dean's Honor List (2021-2022, 2022-2023),  % \section{References}
% References available upon request
\section{Technical Skills}
\begin{tabular}{ l l }
  \textbf{Programming Languages} & C, C++, Python, CUDA, SYCL (DPC++), Bash / POSIX Shell \textit{(+ a lot more, feel free to ask!)}\\
    \textbf{Tooling} & LLVM, MLIR, Linux, Clang (Frontend), Github Actions, GDB/LLDB, Nix/NixOS \\
  %\textbf{Coursework} & Data Structures, Parallel computing, \\
  \textbf{Interests} & (GP)GPU compilers, AI/ML compilers and accelerators, Heterogenous Computing 
    %\textbf{Web Development} & React, React Native, .NET Core, SQL, Express, SASS, Bootstrap, MongoDB, Flask \\
    %\textbf{Certificates} & Microsoft Azure AZ-900, Microsoft Azure AI-900
\end{tabular}

\end{document}
